\begin{definitionbox}{Definition --- Closed Subset of a Topological Space}
\begin{definition}[Closed Subset of a Topological Space]
\label{def:topological-closed-subset}
Let \((X,\tau)\) be a topological space, and let \(F\subseteq X\).  The subset
\(F\) is \emph{closed in \((X,\tau)\)} if its complement
\[
  X\setminus F
\]
is open in \((X,\tau)\).
\end{definition}
\end{definitionbox}

\begin{remark*}[Standard quantified statement]
\[
F\text{ is closed in }(X,\tau)
\Longleftrightarrow
F\subseteq X
\text{ and }
X\setminus F\in\tau.
\]
\end{remark*}

\begin{remark*}[Interpretation]
Closedness is defined by openness of the complement.  A subset is closed when
everything outside it forms an open subset of the ambient topological space.
\end{remark*}

\begin{remark*}[Examples]
The empty set \(\emptyset\) and the whole space \(X\) are both open and
closed in every topological space \((X,\tau)\).  The set \(X\) is open by the
topology axioms and closed because \(X\setminus X=\emptyset\) is open.  The
set \(\emptyset\) is open by the topology axioms and closed because
\(X\setminus\emptyset=X\) is open.
\end{remark*}

\begin{dependencies}
\begin{itemize}
  \item \hyperref[def:topological-space]{Definition: Topological Space}
\end{itemize}
\end{dependencies}

\begin{theorembox}{Theorem}
\begin{theorem}[Closed-Set Characterization of Topologies]
\label{thm:closed-set-characterization-topologies}
Let \(X\) be a set, and let \(\mathcal{C}\) be a collection of subsets of
\(X\).  Suppose that:
\begin{enumerate}
  \item \(\emptyset\in\mathcal{C}\) and \(X\in\mathcal{C}\);
  \item the intersection of any family of elements of \(\mathcal{C}\) belongs
        to \(\mathcal{C}\);
  \item the union of any two elements of \(\mathcal{C}\) belongs to
        \(\mathcal{C}\).
\end{enumerate}
Then there is a unique topology \(\tau\) on \(X\) whose family of closed sets
is exactly \(\mathcal{C}\).

\smallskip
\noindent
\hyperref[prf:closed-set-characterization-topologies]{\textit{Go to proof.}}
\end{theorem}
\end{theorembox}

\begin{remark*}[Standard quantified statement]
\[
\begin{aligned}
&\emptyset,X\in\mathcal{C}
\land
\forall \mathcal{A}\subseteq\mathcal{C}{:}\;
\bigcap_{A\in\mathcal{A}}A\in\mathcal{C}
\\
&{}\land
\forall C_1,C_2\in\mathcal{C}{:}\;
C_1\cup C_2\in\mathcal{C}
\\
&\Longrightarrow
\exists!\tau{:}\;
\operatorname{IsTopology}(X,\tau)
\land
\mathcal{C}
=
\{F\subseteq X:X\setminus F\in\tau\}.
\end{aligned}
\]
\end{remark*}

\begin{remark*}[Predicate reading]
\[
\operatorname{IsTopology}(X,\tau)
\]
means that \(\tau\) is a topology on the carrier set \(X\).
\end{remark*}

\begin{remark*}[Interpretation]
The topology axioms can be stated dually in terms of closed sets.  Open sets
are closed under arbitrary unions and finite intersections, while closed sets
are closed under arbitrary intersections and finite unions.
\end{remark*}

\begin{dependencies}
\begin{itemize}
  \item \hyperref[def:topological-space]{Definition: Topological Space}
  \item \hyperref[def:topological-closed-subset]{Definition: Closed Subset of a Topological Space}
\end{itemize}
\end{dependencies}

\begin{definitionbox}{Definition --- Finer and Coarser Topologies}
\begin{definition}[Finer and Coarser Topologies]
\label{def:finer-coarser-topologies}
Let \(\tau_1\) and \(\tau_2\) be topologies on a set \(X\).  The topology
\(\tau_1\) is \emph{finer than} \(\tau_2\) if
\[
  \tau_2\subseteq\tau_1.
\]
Equivalently, \(\tau_2\) is \emph{coarser than} \(\tau_1\).  The two
topologies are \emph{comparable} if one is finer than the other.
\end{definition}
\end{definitionbox}

\begin{remark*}[Standard quantified statement]
\[
\tau_1\text{ is finer than }\tau_2
\Longleftrightarrow
\tau_2\subseteq\tau_1.
\]
\[
\tau_1\text{ and }\tau_2\text{ are comparable}
\Longleftrightarrow
\tau_2\subseteq\tau_1
\text{ or }
\tau_1\subseteq\tau_2.
\]
\end{remark*}

\begin{remark*}[Interpretation]
A finer topology has at least as many open sets as the coarser topology.
Thus it can distinguish points and subsets at least as sharply as the coarser
topology.
\end{remark*}

\begin{remark*}[Examples]
If \(\mathcal{F}_i\) is the family of closed subsets determined by
\(\tau_i\), then
\[
\tau_1\text{ is finer than }\tau_2
\Longleftrightarrow
\mathcal{F}_2\subseteq\mathcal{F}_1.
\]
The discrete topology is the finest topology on a given set because every
subset is open.
\end{remark*}

\begin{dependencies}
\begin{itemize}
  \item \hyperref[def:topological-space]{Definition: Topological Space}
  \item \hyperref[def:discrete-topology]{Definition: Discrete Topology}
  \item \hyperref[def:topological-closed-subset]{Definition: Closed Subset of a Topological Space}
\end{itemize}
\end{dependencies}
